% ECONOMICS_PROBLEM: {"id": "C73-8", "title": "Payoff functionals guaranteeing a value under one-sided information", "jel_code": "C73", "source_id": "OP-0167"}
% !TeX program = xelatex
\documentclass[11pt,a4paper]{article}
\usepackage[margin=23mm]{geometry}
\usepackage{fontspec}
\setmainfont{Latin Modern Roman}
\usepackage{amsmath,amssymb,mathtools}
\usepackage{xcolor,enumitem,booktabs,longtable,array,xurl,hyperref}
\definecolor{muted}{HTML}{53636C}
\hypersetup{colorlinks=true,urlcolor=blue,linkcolor=blue}
\setlength{\parindent}{0pt}
\setlength{\parskip}{5pt}
\newcommand{\R}{\mathbb R}
\newcommand{\E}{\mathbb E}
\newcommand{\N}{\mathbb N}
\newcommand{\ind}{\mathbf 1}
\newcommand{\Var}{\operatorname{Var}}
\newcommand{\Cov}{\operatorname{Cov}}
\newcommand{\supp}{\operatorname{supp}}
\DeclareMathOperator*{\argmin}{arg\,min}
\DeclareMathOperator*{\argmax}{arg\,max}
\newcommand{\entrymeta}[3]{{\sffamily\small\color{muted}\textbf{\texttt{#1}}\quad|\quad #2\par #3\par}}
\newcommand{\Problemref}[1]{\texttt{#1}}
\newcommand{\JELref}[2]{\texttt{#2} (JEL \texttt{#1})}
\begin{document}
\section*{Payoff functionals guaranteeing a value under one-sided information}
\textbf{Problem ID:} \texttt{C73-8}\quad \textbf{JEL:} \texttt{C73}\par
% BEGIN ECONOMICS STATEMENT
\label{problem:OP-0167}
% GAME_THEORY_PROBLEM: {"local_id": "GT-037", "original_number": 37, "kind": "P", "entry_kind": "statement_only", "published": false, "review_required": true, "category": "game_theory"}
\label{game:GT-037}
{\sffamily\small\color{muted}
{\sffamily\small\color{muted}\textbf{GT-037} | Repeated games with incomplete information\par P | Source-reported characterization question\par}
\par}
\subsection*{Definitions and mathematical statement}
Let $K,I,J$ be nonempty finite sets. Nature chooses $k\in K$ according to $p\in\Delta(K)$ and reveals it only to the maximizing player. At every date the two players simultaneously choose actions in $I$ and $J$ and then both observe the pair. Write $\mathcal H=\bigcup_{t\geq0}(I\times J)^t$. The informed player's behavioral strategy is $\sigma:K\times\mathcal H\to\Delta(I)$ and the uninformed player's is $\tau:\mathcal H\to\Delta(J)$. For a bounded Borel functional $f:K\times(I\times J)^{\mathbb N}\to\mathbb R$, set
\[
 \gamma_{p,f}(\sigma,\tau)=\mathbb E_{p,\sigma,\tau}f(k,(i_t,j_t)_{t\geq0}),
\]
\[
 \underline v_f(p)=\sup_\sigma\inf_\tau\gamma_{p,f}(\sigma,\tau),
 \qquad
 \overline v_f(p)=\inf_\tau\sup_\sigma\gamma_{p,f}(\sigma,\tau).
\]
For fixed $K,I,J$, characterize intrinsically the class
\[
 \mathcal V_{K,I,J}=\{f\text{ bounded Borel}:
 \forall p\in\Delta(K),\ \underline v_f(p)=\overline v_f(p)\}.
\]
A characterization should give conditions on $f$ and the information/action structure, rather than merely restating equality of the two optimization problems.
\subsection*{Short English statement}
Which bounded infinite-play payoffs ensure that the informed and uninformed players have the same minimax value?
\subsection*{Variants and scope boundaries}
Membership at one specified prior is weaker than prior-uniform membership. Tail measurability means that changing finitely many action coordinates leaves $f(k,\cdot)$ unchanged; shift invariance means invariance under deletion of an initial action pair. These notions differ. Borel measurability, complete action monitoring, and finite state/action alphabets are fixed here.
\subsection*{Source relationship and status}
[S1], Section 5, asks for payoff properties guaranteeing a value. Its tail-measurable maxmin theorem does not assert that the upper and lower values coincide; Section 4 provides relevant nonexistence examples. The prior-uniform class above is an explicit organizational choice.
\subsection*{References}
\begingroup\footnotesize\raggedright
\textbf{[S1]} Gil Bar Castellon Koltun, Ehud Lehrer and Eilon Solan, \emph{The Maxmin Value of Repeated Games with Incomplete Information on One Side and Tail-Measurable Payoffs}, arXiv:2512.01581 (2025). Locators: Definitions 2.1--2.4, pp. 6--7; Section 4 (examples); Section 5, p. 18 (characterization questions). \url{https://arxiv.org/pdf/2512.01581}
\par\endgroup

\subsection*{Common conventions: Source mathematical conventions}

\textbf{Finite games.} For a finite set $A$, $\Delta(A)$ is its probability
simplex. Players choose independent mixed strategies in Nash equilibrium;
correlation is allowed only when explicitly part of the solution concept.
In a game with payoff functions $u_i$ and strategy profile $x$, an additive
$\varepsilon$ Nash equilibrium satisfies
\[
 u_i(x)\geq u_i(x_i',x_{-i})-\varepsilon
 \quad\text{for every player $i$ and every admissible deviation $x_i'$}.
\]
For numerical approximation, payoffs are normalized as locally specified.
An $\varepsilon$ payoff residual is distinct from distance at most
$\varepsilon$ to the set of exact equilibria. Point convergence, convergence
of empirical play, and convergence in distance to an equilibrium set are
also distinct.

\textbf{Computational representations.} Unless stated otherwise, a complexity
question takes rational data with binary encoding; polynomial time is measured
in total bit length. A fixed-accuracy polynomial algorithm is not necessarily
polynomial in $1/\varepsilon$, and neither is necessarily polynomial in
$\log(1/\varepsilon)$. Succinct games and value, demand, or sampling oracles
must declare their interfaces. Exact real solutions require an explicit
representation; an unspecified list of arbitrary real numbers is not a
finite computational output. A claimed algorithmic barrier is meaningful
only for its specified model and reductions.

\textbf{Dynamic games.} Histories, information, observation, and allowable
randomization are part of the model. A stationary strategy depends only on
the current state; a positional strategy in a deterministic graph game
chooses one action at each controlled vertex. Nash equilibrium at an initial
state and equilibrium after every history are different requirements.
An expectation of a pathwise $\liminf$ payoff, a $\liminf$ of expected averages,
a limiting expected average, and a uniform finite-horizon equilibrium are
different criteria. None is silently substituted for another.

\textbf{Allocation and mechanisms.} A complete allocation assigns every item
exactly once unless otherwise stated. Zero-valued goods, negative values,
capacity constraints, feasibility, lotteries, and copies of goods affect
fairness definitions and are specified locally. Dominant-strategy incentive
compatibility, Bayesian incentive compatibility, universal truthfulness,
and truthfulness in expectation impose different inequalities. Whenever
an objective ratio has zero denominator, its convention is stated or those
instances are excluded.

\textbf{Infinite-dimensional models.} Expectations, integrals, stochastic
equations, and optimization objectives require their local regularity,
integrability, filtrations, and admissible domains. A backward--forward
mean-field system is a boundary-value problem; it is not an initial-value
semiflow without an additional construction. Long-time, large-population,
and vanishing-noise limits require an order of limits and a convergence
topology. If a minimum need not be attained, the objective is its infimum
or supremum and the approximation requirement is stated separately.
% END ECONOMICS STATEMENT
\end{document}
